% page 283 du fac-simile (image p-282 du scan)
% bloc 162.6 x 246.3 mm ; marge gauche 39.4 mm
\pgc{17.780mm}{0.000mm}{
\l{}
\l{}
\l{}
\l{\cel{55}275}
\l{}
\l{\sou{kemio}.\cel{2}La cienco qua studias la konsisto intima di la korpi}
\l{\cel{3}diversa (simpla o komplexa), la +leyi (legi) segun qui inter-}
\l{\cel{3}kombinas la elementi de qui oli naskas, ed olia qualesi}
\l{\cel{3}propra. - DEFIRS.}
\l{}
\l{\sou{kemiotaktismo}. (\sou{biol.}) Atrakto o repulso, che enti vivanta,}
\l{\cel{3}da substanco kemiala.}
\l{}
\l{\sou{kemiotropismo}.\cel{2}(\sou{biol.}) Orientizo, direkto di la planti, per}
\l{\cel{3}reaktivi kemiala.}
\l{}
\l{\sou{kenelo.}\cel{2}Buleto oblonga de karno, de fisho-karno, uzata kom}
\l{\cel{3}disho-garnituro. - EF.}
\l{}
\l{\sou{kepio}. I. Kapo-vesto lejera, kun mikra viziero, quan +weras}
\l{\cel{3}ordinare la armeani. - II. Kapo-vesto sama-forma di la skol-}
\l{\cel{3}ani, kolegiani, liceani, e c. - DEFIRS.}
\l{}
\l{\sou{keratino}. Substanco albuminoida, tre sulfoza, qua esas elemento}
\l{\cel{3}di la tisui sintezala (tegumenti, faneri). - SF.}
\l{}
\l{\sou{kerlo}. Viro vivaca e joyoza. - D.}
\l{}
\l{\sou{kermeso}. (\sou{zool.}) Kochenilo qua vivas sur speco di querko, e}
\l{\cel{3}di qua la femino facas,sur la folii, la stipi e la branchi}
\l{\cel{3}di ta arboro, mikra sheli sferatra e reda. - L.\cel{1}\sou{coccus ilicia}}
\l{\cel{3}- DEFIRS.}
\l{}
\l{kerno. Parto harda, ligna, qua esas interne di ula frukti,}
\l{\cel{3}e qua kontenas la \textquotedbl{}mandelo\textquotedbl{}, la grano. - DE.}
\l{}
\l{\sou{kerozeno}. (\sou{kemio}).Lum-oleo, obtenita per distilar petrol-olei}
\l{\cel{3}kruda. - EF.}
\l{}
\l{\sou{kerub}o.I. En la Testamento Olda (biblo) : reprezenturo de la an-}
\l{\cel{3}jeli, qui ornis la tabernaklo ; nomo di ula anjeli. - II.}
\l{\cel{3}Che la kristani : anjelo qua, en la hierarkio cielala, esas}
\l{\cel{3}situita pos la arkianjelo, e quan la pikturi e skulturi re-}
\l{\cel{3}prezentas ordinare ye kapo di infanto aloza. - DEFIRS.}
\l{}
\l{\sou{kesto}. Quaza buxo por la transporto ; asemblajo ek planki (ula-}
\l{\cel{3}kaze : ajuroza, kun fundo e kovrilo fixigita per klovi).}
\l{}
\l{\textquotedbl{}\sou{khan\textquotedbl{}}. I. Sinioro (che la hordi nomada, Turka o Mongola).}
\l{\cel{3}- II.Suvereno (olim, che la Turki e la Tatari.)\cel{2}- III.}
\l{\cel{3}(en\cel{1}\sou{Persi}a)Guvernisto di provinco, ed ulagrade oficiero}
\l{\cel{3}statala.}
\l{}
\l{\textquotedbl{}\sou{khedivo\textquotedbl{}}.\cel{2}La suvereno di Egiptia,olim.}
\l{}
\l{\sou{kikar}. (\sou{netrans}.)(\sou{aludante kavalo, asno, e c.)}\cel{2}Lansar rapide}
\l{\cel{3}addope sua membri dopa,apogante su sur la membri avana. - E.}
}
